\documentclass{amsart}
\usepackage{amsmath,amssymb,amsthm,hyperref}
\newtheorem{theorem}{Theorem}
\newtheorem{lemma}{Lemma}
\newtheorem{proposition}{Proposition}
\theoremstyle{definition}
\newtheorem{construction}{Construction}
\begin{document}

\title{A $K$-metric making $A$ a contraction}
\maketitle

\noindent\textbf{Standing hypotheses.} Let $X$ be a nonempty set, $A:X\to X$ an
operator, $A^0=\mathrm{id}_X$, $A^{n+1}=A\circ A^n$, and
$F_{A^n}=\{x\in X: A^n(x)=x\}$. We assume there is $x^*\in X$ with
\begin{equation}\label{eq:hyp}
  F_A=F_{A^n}=\{x^*\}\qquad\text{for every }n\ge 1 .
\end{equation}

We will in fact produce an \emph{ordinary} (real-valued) metric for which $A$ is a
Banach contraction; this is a special case of a $K$-metric contraction, obtained
by taking $K=\mathbb{R}$ and $Q=\lambda\,\mathrm{id}_{\mathbb{R}}$ for a constant
$\lambda\in(0,1)$. We state the choice of $K$, $Q$ and $d$ first, and then prove
that all required properties hold.

\bigskip
\noindent\textbf{The construction.}
\begin{itemize}
\item \emph{Ordered Banach space.} $K=\mathbb{R}$ with the usual norm $\lVert
t\rVert=|t|$ and positive cone $\mathbb{R}_{\ge 0}=[0,\infty)$; thus $t\ge 0$ in
$K$ means $t\ge 0$ in the usual sense.
\item \emph{Contraction operator.} Fix once and for all a number
$\lambda\in(0,1)$ (e.g.\ $\lambda=\tfrac12$). Let $Q:\mathbb{R}\to\mathbb{R}$,
$Q(t)=\lambda t$. Then $Q$ is linear, positive (it maps $[0,\infty)$ into
$[0,\infty)$), and $\lVert Q\rVert=\lambda<1$.
\item \emph{The $K$-metric.} We construct (below) a function
$\varphi:X\to[0,\infty)$ with
\begin{equation}\label{eq:phi}
\varphi(x)=0\iff x=x^*,\qquad
\varphi\bigl(A(x)\bigr)\le\lambda\,\varphi(x)\ \ \text{for all }x\in X,
\end{equation}
and then set
\begin{equation}\label{eq:d}
  d(x,y)=
  \begin{cases}
    0, & x=y,\\
    \varphi(x)+\varphi(y), & x\ne y.
  \end{cases}
\end{equation}
\end{itemize}

The whole problem thus reduces to constructing $\varphi$ as in \eqref{eq:phi}, and
to checking that \eqref{eq:d} is a $K$-metric for which $A$ is a contraction with
operator $Q$. We do the verification first (Sections~1--2), and then carry out the
construction of $\varphi$ (Section~3).

\section{The map \eqref{eq:d} is a $K$-metric}

Assume $\varphi$ satisfies \eqref{eq:phi}. We check axioms (i)--(iii) for $d$ in
\eqref{eq:d}.

\smallskip
\noindent\emph{(i) Positivity and faithfulness.} Each value $d(x,y)$ is a sum of
nonnegative reals, hence $d(x,y)\ge 0$ in $K$. If $x=y$ then $d(x,y)=0$ by
definition. Conversely, suppose $x\ne y$; then
$d(x,y)=\varphi(x)+\varphi(y)$. Since $x\ne y$, at least one of $x,y$ differs from
$x^*$, say $x\ne x^*$; then $\varphi(x)>0$ by \eqref{eq:phi}, and
$\varphi(y)\ge 0$, so $d(x,y)>0$. Hence $d(x,y)=0\iff x=y$.

\smallskip
\noindent\emph{(ii) Symmetry.} Immediate from the symmetry of \eqref{eq:d} in
$x,y$.

\smallskip
\noindent\emph{(iii) Triangle inequality.} Fix $x,y,z\in X$; we show
$d(x,z)\le d(x,y)+d(y,z)$.

\emph{Case $x=z$.} Then $d(x,z)=0\le d(x,y)+d(y,z)$, because the right-hand side
is a sum of nonnegative reals.

\emph{Case $x\ne z$.} Then $d(x,z)=\varphi(x)+\varphi(z)$. We distinguish
according to $y$:
\begin{itemize}
\item If $y=x$, then $d(x,y)=0$ and (since $y\ne z$ because $x\ne z$)
$d(y,z)=\varphi(y)+\varphi(z)=\varphi(x)+\varphi(z)=d(x,z)$, so equality holds.
\item If $y=z$, symmetrically $d(x,y)=\varphi(x)+\varphi(z)=d(x,z)$ and
$d(y,z)=0$, so equality holds.
\item If $y\notin\{x,z\}$, then
\[
  d(x,y)+d(y,z)=\bigl(\varphi(x)+\varphi(y)\bigr)+\bigl(\varphi(y)+\varphi(z)\bigr)
  =\varphi(x)+\varphi(z)+2\varphi(y)\ge\varphi(x)+\varphi(z)=d(x,z),
\]
using $\varphi(y)\ge 0$.
\end{itemize}
In all cases $d(x,z)\le d(x,y)+d(y,z)$.

Thus $d$ is a $K$-metric on $X$.

\section{$A$ is a $K$-metric contraction}

We show $d\bigl(A(x),A(y)\bigr)\le Q\bigl(d(x,y)\bigr)=\lambda\,d(x,y)$ for all
$x,y\in X$.

\emph{Case $A(x)=A(y)$.} Then $d\bigl(A(x),A(y)\bigr)=0\le\lambda\,d(x,y)$,
because $d(x,y)\ge 0$ and $\lambda>0$.

\emph{Case $A(x)\ne A(y)$.} Then necessarily $x\ne y$, so
$d(x,y)=\varphi(x)+\varphi(y)$, and
\[
  d\bigl(A(x),A(y)\bigr)=\varphi\bigl(A(x)\bigr)+\varphi\bigl(A(y)\bigr)
  \le\lambda\,\varphi(x)+\lambda\,\varphi(y)
  =\lambda\bigl(\varphi(x)+\varphi(y)\bigr)=\lambda\,d(x,y),
\]
where the inequality uses the second part of \eqref{eq:phi} applied to $x$ and to
$y$.

Hence $d\bigl(A(x),A(y)\bigr)\le Q\bigl(d(x,y)\bigr)$ for all $x,y\in X$, i.e.\ $A$
is a $K$-metric contraction on $(X,d)$ with operator $Q$ of norm $\lambda<1$.

\section{Construction of $\varphi$ satisfying \eqref{eq:phi}}

It remains to construct $\varphi:X\to[0,\infty)$ with $\varphi(x)=0\iff x=x^*$ and
$\varphi(A(x))\le\lambda\,\varphi(x)$ for all $x$. Throughout we use only
hypothesis \eqref{eq:hyp}; its sole consequence we need is the following.

\begin{lemma}[No nontrivial periodic points]\label{lem:nocycle}
If $z\in X$ and $A^k(z)=z$ for some $k\ge 1$, then $z=x^*$.
\end{lemma}
\begin{proof}
$A^k(z)=z$ means $z\in F_{A^k}$, and $F_{A^k}=\{x^*\}$ by \eqref{eq:hyp}; hence
$z=x^*$.
\end{proof}

\begin{lemma}[Orbit dichotomy]\label{lem:dich}
Let $z\in X$. Exactly one of the following holds:
\begin{enumerate}
\item[\rm(a)] $A^m(z)=x^*$ for some $m\ge 0$ (we say $z$ is in the \emph{basin}
$B$); or
\item[\rm(b)] the iterates $z,A(z),A^2(z),\dots$ are pairwise distinct and none of
them equals $x^*$ (we say $z\in X':=X\setminus B$).
\end{enumerate}
\end{lemma}
\begin{proof}
If (a) fails, then $A^m(z)\ne x^*$ for all $m\ge 0$. Suppose two iterates
coincided, say $A^i(z)=A^j(z)$ with $0\le i<j$. Put $w:=A^i(z)$. Then
$A^{j-i}(w)=A^{j}(z)=A^i(z)=w$ with $j-i\ge 1$, so $w=x^*$ by
Lemma~\ref{lem:nocycle}; but $w=A^i(z)\ne x^*$, a contradiction. Hence all
iterates are distinct, and none equals $x^*$; this is (b). Clearly (a) and (b) are
mutually exclusive.
\end{proof}

For $z\in B$ define
\begin{equation}\label{eq:Ndef}
  N(z):=\min\{\,n\ge 0: A^n(z)=x^*\,\}\in\{0,1,2,\dots\},
\end{equation}
the (finite) entrance time into $x^*$. Note $N(x^*)=0$.

\begin{lemma}[Forward invariance of $B$ and $X'$; behaviour of $N$]\label{lem:inv}
\hfill
\begin{enumerate}
\item[\rm(a)] If $z\in B$ then $A(z)\in B$, and if moreover $z\ne x^*$ then
$N\bigl(A(z)\bigr)=N(z)-1$.
\item[\rm(b)] If $z\in X'$ then $A(z)\in X'$.
\end{enumerate}
\end{lemma}
\begin{proof}
(a) If $A^m(z)=x^*$ then $A^{m}(z)=x^*$ shows $z\in B$, and $A^{m-1}(A(z))=x^*$
when $m\ge 1$, so $A(z)\in B$ as well; if $m=0$ then $z=x^*$ and
$A(z)=x^*\in B$. Now let $z\in B$, $z\ne x^*$, so $N(z)=:k\ge 1$ and $A^k(z)=x^*$,
$A^{k-1}(z)\ne x^*$. Then $A^{k-1}(A(z))=A^{k}(z)=x^*$, so
$N(A(z))\le k-1$. If $N(A(z))=:\ell<k-1$, then
$A^{\ell+1}(z)=A^{\ell}(A(z))=x^*$ with $\ell+1<k$, contradicting minimality of
$N(z)=k$. Hence $N(A(z))=k-1=N(z)-1$.

(b) If $A(z)\in B$, then $A^m(A(z))=x^*$ for some $m\ge 0$, i.e.\
$A^{m+1}(z)=x^*$, so $z\in B$. By contraposition, $z\in X'$ implies $A(z)\in X'$.
\end{proof}

\subsection*{Definition of $\varphi$ on the basin $B$}
For $z\in B$ set
\begin{equation}\label{eq:phiB}
  \varphi(z):=\lambda^{-N(z)}-1 .
\end{equation}
Since $0<\lambda<1$ we have $\lambda^{-N(z)}\ge 1$, so $\varphi(z)\ge 0$.
Moreover $\varphi(z)=0\iff\lambda^{-N(z)}=1\iff N(z)=0\iff z=x^*$ (the last
equivalence holds because $N(z)=0$ means $A^0(z)=z=x^*$). Thus on $B$ the value
$\varphi$ is $0$ exactly at $x^*$ and positive elsewhere.

We verify $\varphi(A(z))\le\lambda\,\varphi(z)$ for $z\in B$.
If $z=x^*$, then $A(z)=x^*$ and $\varphi(A(z))=0=\lambda\cdot 0=\lambda\varphi(z)$.
If $z\ne x^*$, write $k:=N(z)\ge 1$; by Lemma~\ref{lem:inv}(a),
$N(A(z))=k-1$, so
\[
  \varphi\bigl(A(z)\bigr)=\lambda^{-(k-1)}-1=\lambda^{-k+1}-1,
  \qquad
  \lambda\,\varphi(z)=\lambda\bigl(\lambda^{-k}-1\bigr)=\lambda^{-k+1}-\lambda .
\]
Since $\lambda<1$ we have $-1<-\lambda$, hence
$\varphi(A(z))=\lambda^{-k+1}-1<\lambda^{-k+1}-\lambda=\lambda\,\varphi(z)$. Thus
$\varphi(A(z))\le\lambda\,\varphi(z)$ on $B$, as required.

\subsection*{Definition of $\varphi$ on $X'$}
On $X'$ we use the following relation. For $x,y\in X'$ write $x\sim y$ if there
exist $m,n\ge 0$ with $A^m(x)=A^n(y)$ (eventual coincidence of forward orbits).

\begin{lemma}\label{lem:equiv}
$\sim$ is an equivalence relation on $X'$, and each $\sim$-class is contained in
$X'$.
\end{lemma}
\begin{proof}
Reflexivity ($m=n=0$) and symmetry are clear. For transitivity, let
$A^{m}(x)=A^{n}(y)$ and $A^{p}(y)=A^{q}(z)$. Applying $A^{p}$ to the first and
$A^{n}$ to the second gives $A^{m+p}(x)=A^{n+p}(y)=A^{n+q}(z)$, so $x\sim z$.
(Membership in $X'$ holds by definition of the relation on $X'$.)
\end{proof}

\begin{lemma}[Well-defined relative level]\label{lem:delta}
Let $x,y\in X'$ with $x\sim y$. If $A^m(x)=A^n(y)$ and $A^{m'}(x)=A^{n'}(y)$ with
$m,n,m',n'\ge 0$, then $m-n=m'-n'$.
\end{lemma}
\begin{proof}
Without loss of generality $m'\ge m$. Apply $A^{m'-m}$ to $A^m(x)=A^n(y)$:
\[
  A^{m'}(x)=A^{\,n+m'-m}(y).
\]
Comparing with $A^{m'}(x)=A^{n'}(y)$ yields
$A^{n'}(y)=A^{\,n+m'-m}(y)$. Set $p:=n'$ and $q:=n+m'-m$; both are $\ge 0$ (indeed
$q\ge n\ge0$ since $m'\ge m$). If $p\ne q$, say $p<q$, put $w:=A^{p}(y)\in X'$ (by
Lemma~\ref{lem:inv}(b), $X'$ is forward invariant); then $A^{q-p}(w)=w$ with
$q-p\ge 1$, so $w=x^*$ by Lemma~\ref{lem:nocycle}, contradicting $w\in X'$. Hence
$p=q$, i.e.\ $n'=n+m'-m$, i.e.\ $m-n=m'-n'$.
\end{proof}

By Lemma~\ref{lem:delta}, for $x\sim y$ in $X'$ the integer
\begin{equation}\label{eq:rel}
  \mathrm{rel}(x,y):=m-n,\qquad\text{where }A^m(x)=A^n(y),
\end{equation}
is well defined (independent of the choice of $m,n$).

Now, using the Axiom of Choice, pick a representative $r_C\in C$ in each
$\sim$-equivalence class $C\subseteq X'$. For $x\in X'$ lying in class $C$ define
\begin{equation}\label{eq:phiXp}
  \varphi(x):=\lambda^{-\,\mathrm{rel}(x,r_C)} .
\end{equation}
Since $\lambda>0$, we have $\varphi(x)>0$ for every $x\in X'$ (note $x^*\notin
X'$, so positivity here is consistent with the requirement $\varphi=0$ only at
$x^*$).

We verify $\varphi(A(x))=\lambda\,\varphi(x)$ for $x\in X'$ (the equality is
stronger than the required inequality). Let $x\in X'$ belong to class $C$ with
representative $r:=r_C$, and choose $m,n\ge 0$ with $A^m(x)=A^n(r)$, so
$\mathrm{rel}(x,r)=m-n$. Then $A(x)\in X'$ (Lemma~\ref{lem:inv}(b)), and $A(x)\sim
x\sim r$, so $A(x)$ lies in the same class $C$. Applying $A^m$ to $A(x)$:
\[
  A^{m}\bigl(A(x)\bigr)=A^{m+1}(x)=A\bigl(A^m(x)\bigr)=A\bigl(A^n(r)\bigr)=A^{n+1}(r),
\]
so $(m,n+1)$ is an admissible pair for $A(x)$ and $r$, giving
$\mathrm{rel}(A(x),r)=m-(n+1)=\mathrm{rel}(x,r)-1$. Therefore
\[
  \varphi\bigl(A(x)\bigr)=\lambda^{-(\mathrm{rel}(x,r)-1)}
  =\lambda\cdot\lambda^{-\mathrm{rel}(x,r)}=\lambda\,\varphi(x)
  \le\lambda\,\varphi(x).
\]

\subsection*{Conclusion of the construction}
The two definitions \eqref{eq:phiB} on $B$ and \eqref{eq:phiXp} on
$X'=X\setminus B$ give a single function $\varphi:X\to[0,\infty)$, because $B$ and
$X'$ partition $X$ (Lemma~\ref{lem:dich}). We have shown:
\begin{itemize}
\item $\varphi(z)\ge 0$ for all $z$, with $\varphi(z)=0$ iff $z=x^*$ (positive on
$B\setminus\{x^*\}$ and on $X'$; zero at $x^*$);
\item $\varphi(A(z))\le\lambda\,\varphi(z)$ for all $z$ (with equality on $X'$ and
the verified inequality on $B$).
\end{itemize}
This is exactly \eqref{eq:phi}. We must finally check that the relation
$\varphi(A(z))\le\lambda\varphi(z)$ is internally consistent across the partition,
i.e.\ that for every $z\in X$ both $z$ and $A(z)$ are evaluated by formulas already
established. Indeed $B$ and $X'$ are each forward invariant
(Lemma~\ref{lem:inv}), so if $z\in B$ then $A(z)\in B$ and the inequality was
checked via \eqref{eq:phiB}; if $z\in X'$ then $A(z)\in X'$ and the equality was
checked via \eqref{eq:phiXp}. Hence \eqref{eq:phi} holds globally.

\section{Result}

\begin{theorem}
Let $X$ be a nonempty set and $A:X\to X$ satisfy \eqref{eq:hyp}. Fix any
$\lambda\in(0,1)$. With $K=\mathbb{R}$ (positive cone $[0,\infty)$, usual norm),
$Q=\lambda\,\mathrm{id}_{\mathbb{R}}$ (positive, linear, $\lVert
Q\rVert=\lambda<1$), $\varphi$ constructed in \eqref{eq:phiB}--\eqref{eq:phiXp},
and $d$ defined by \eqref{eq:d}, the map $d$ is a $K$-metric on $X$ and
\[
  d\bigl(A(x),A(y)\bigr)\le Q\bigl(d(x,y)\bigr)\qquad\text{for all }x,y\in X.
\]
Thus $A:(X,d)\to(X,d)$ is a $K$-metric contraction.
\end{theorem}
\begin{proof}
Section~3 produces $\varphi$ satisfying \eqref{eq:phi}. Given \eqref{eq:phi},
Section~1 shows $d$ is a $K$-metric and Section~2 shows the contraction inequality
with $Q(t)=\lambda t$, which is positive, linear, and has $\lVert
Q\rVert=\lambda<1$.
\end{proof}

\bigskip
\noindent\textbf{Remark.} The construction yields a genuine real-valued metric
(the case $K=\mathbb{R}$, $Q=\lambda\,\mathrm{id}$), so in particular $A$ is an
ordinary Banach $\lambda$-contraction for this metric. This is precisely Bessaga's
converse to the Banach contraction principle, here phrased and proved within the
$K$-metric framework requested by the problem. Only hypothesis \eqref{eq:hyp} was
used, through its single consequence Lemma~\ref{lem:nocycle} (absence of
nontrivial periodic points).

\end{document}
